\Needspace{10\baselineskip}
\centerline{\textsc{Gauss to Bode.}}

\medskip
Astronomisches Jahrbuch für 1808. Pp. 187---190. Berlin 1808.

\medskip
\begin{flushright}
Braunschweig, 30 September 1804.
\end{flushright}

I take the liberty of presenting here to Your --- all my observations of Juno to date. I used only a ring micrometer, and took as much care as possible in making the observations. They are all the result of very numerous comparisons, some with stars from \textsc{Piazzi}’s catalogue and some with stars from the Hist. Céleste, which I carefully reduced according to that catalogue. With the exception of the first 3, which were made with an inferior instrument, all were made with a very good \textsc{Short} reflecting telescope.

\begin{center}
\begingroup\fontsize{10}{12}\selectfont\setlength{\tabcolsep}{3pt}\renewcommand{\arraystretch}{1.12}
\begin{tabular}{r|rrr|rrr|rrr}
1804 & \multicolumn{3}{c|}{Mean time} & \multicolumn{6}{c}{Apparent} \\ \cline{5-10}
 & \multicolumn{3}{c|}{} & \multicolumn{3}{c|}{Right ascension} & \multicolumn{3}{c}{Southern declination} \\ \hline
12 Sept. & $10^u$ & $35^m$ & $2^s$ & $0^\circ$ & $54'$ & $31''$ & $1^\circ$ & $38'$ & $15''$ \\
13 & $9$ & $41$ & $32$ &  & $45$ & $24$ &  & $50$ & $59$ \\
14 & $11$ & $38$ & $49$ &  & $35$ & $37$ & $2$ & $5$ & $1$ \\
15 & $10$ & $16$ & $17$ &  & $26$ & $53$ &  & $17$ & $35$ \\
16 & $10$ & $37$ & $4$ &  & $17$ & $17$ &  & $31$ & $20$ \\
17 & $11$ & $28$ & $59$ &  & $7$ & $23$ &  & $44$ & $29$ \\
\end{tabular}\endgroup
\end{center}
\begin{center}
\begingroup\fontsize{10}{12}\selectfont\setlength{\tabcolsep}{3pt}\renewcommand{\arraystretch}{1.12}
\begin{tabular}{r|rrr|rrr|rrr}
1804 & \multicolumn{3}{c|}{Mean time} & \multicolumn{6}{c}{Apparent} \\ \cline{5-10}
 & \multicolumn{3}{c|}{} & \multicolumn{3}{c|}{Right ascension} & \multicolumn{3}{c}{Southern declination} \\ \hline
18 Sept. & $11^u$ & $22^m$ & $16^s$ & $359^\circ$ & $57'$ & $26''$ &  & $58'$ & $3''$ \\
21 & $10$ & $24$ & $52$ &  & $28$ & $6$ & $3$ & $38$ & $0$ \\
24 & $10$ & $1$ & $2$ & $358$ & $57$ & $53$ & $4$ & $18$ & $6$ \\
25 & $8$ & $44$ & $25$ &  & $48$ & $12$ &  & $30$ & $44$ \\
27 & $10$ & $20$ & $29$ &  & $27$ & $20$ &  & $57$ & $47$ \\
28 & $8$ & $29$ & $4$ &  & $18$ & $20$ & $5$ & $10$ & $22$ \\
\end{tabular}\endgroup
\end{center}
\SourceLeaf{254}%

Drawing in part on these observations and in part on those of our mutual friend Dr. \textsc{Olbers} and on some from Baron \textsc{von Zach}, I have already attempted a first approximation to the elements, which I have the honour to communicate to you here. I consider Juno’s planetary nature to have been established by now \emph{from the observations}, and I venture to think that the elements already suffice to give some idea of the approximate dimensions of the orbit,

\begin{center}
I\textsuperscript{st} elements of Juno’s orbit.
\end{center}

\begin{center}
\begingroup\fontsize{8}{10}\selectfont\setlength{\tabcolsep}{3pt}
\begin{tabular}{lr}
Epoch, 5 September 1804, $0^u$ mean time on the Seeberg meridian . & $20^\circ\ 38'\ 56''$ \\
Mean daily motion . . . . . . . . . . & $779''80$ \\
Aphelion . . . . . . . . . . . . . . . . & $239^\circ\ 14'\ 2''$ \\
Log. of the semiaxis . . . . . . . . . . . . . & $0.438682$ \\
Eccentricity . . . . . . . . . . . . . . . & $0.287359$ \\
Ascending node . . . . . . . . . . . . . . & $171^\circ\ 15'\ 35''$ \\
Inclination . . . . . . . . . . . . . . . . . . & $13\ 34\ 59$ \\
\end{tabular}\endgroup
\end{center}

Naturally, since the period covered by the observations is still so short, these elements will still need considerable improvement. In particular, the uncertainty in the mean motion is still much greater than the small difference from the mean motion of $\CeresSymbol$ and $\PallasSymbol$, so all 3 may well come much closer to equality.

To do what I can to help astronomers farther away who may still have been unable to find the planet, I have used the elements above to calculate a short ephemeris that will, I hope, show its course with enough accuracy for them to find it over the next several weeks.

\Needspace{26\baselineskip}
\begin{flushright}
Braunschweig, 12 November 1804.
\end{flushright}

In offering Your--- my most sincere thanks for the observations you sent me, I also continue my account of the latest results of my investigations of Juno’s orbit, although they were already calculated 14 days ago and will therefore soon admit of further improvement. Here are the

\begin{center}
III\textsuperscript{rd} elements of Juno.
\end{center}

\begin{center}
\begingroup\fontsize{9}{11}\selectfont\setlength{\tabcolsep}{3pt}
\begin{tabular}{lr}
Epoch 1805 on the Seeberg meridian . . . . . . . . & $43^\circ\ 21'\ 2''$ \\
Aphelion . . . . . . . . . . . . . . . . . & $233\ 56\ 6$ \\
Inclination . . . . . . . . . . . . . . . . . . & $12\ 52\ 48$ \\
Logarithm of the semiaxis . . . . . . . . . . . . & $0.426699$ \\
Daily motion . . . . . . . . . . . . . . & $812''754$ \\
Node . . . . . . . . . . . . . . . . . . & $171^\circ\ 0'\ 0''$ \\
Eccentricity . . . . . . . . . . . . . . . . & $0.263182$ \\
\end{tabular}\endgroup
\end{center}

I already regard it as settled that Juno’s mean motion differs from that of $\CeresSymbol$ and $\PallasSymbol$ in being greater.

When I calculated the III\textsuperscript{rd} elements, I had Mr. \textsc{von Zach}’s observations through 21 October before me, and the elements agree with them within a few seconds.
So far I have Dr. \textsc{Maskelyne}’s observations through the 17\textsuperscript{th}. The later observations, however, establish that these elements now give too large a right ascension. As for the declination, I cannot decide anything yet. Here is the comparison with the latest observations.
\SourceLeaf{255}%


\begin{center}\begin{tabular}{r|cc|c|cc}
 & \multicolumn{2}{c|}{Prof. \textsc{Bode}} & & \multicolumn{2}{c}{Dr. \textsc{Olbers}} \\
 & R.A. & Decl. & & R.A. & Decl. \\
\hline
20 Oct. & $+33''$ & $+7''$ & 27 Oct. & $+21''$ & $+15''$ \\
2 Nov. & $+48$ & $+5$ & 30 & $+34$ & $+12$ \\
5 & $+46$ & $-15$ & 2 Nov.  & $+33$ & $-1$ \\
 & & & 5 & $+39$ & $+25$ \\
 & & & 6 & $+37$ & $+32$ \\
\end{tabular}\end{center}

On 20 December $\CeresSymbol$ and $\JunoSymbol$ will come so close together that it will be possible to observe them simultaneously through the telescope.


